\section{When Next-Edge Cost Is Sufficient}
\label{sec:next}
The next-edge variant keeps the UCB exploration term and truncates only the cost horizon. Write its index and the complete cost decomposition as
\begin{equation}
\begin{split}
I_t^{\mathrm{next}}(a)&=\widehat Q_t(a)+b_t(a)-\lambda e_t(a),\\
c_t(a)&=e_t(a)+r_t(a)+h_t.
\end{split}
\label{eq:next}
\end{equation}
Here $e_t(a)$ is the next-edge cost, $r_t(a)$ is the omitted action-dependent remainder, and $h_t$ is any action-independent cost at this node and time. The common $h_t$ cancels from both rankings and regret gaps. The word ``next'' does not mean greedy reward selection: exploration is retained.

\begin{proposition}[Exact reduction and approximation error]
\label{prop:next}
If $r_t(a)$ is the same for every action, next-edge and complete-cost indices have exactly the same maximizers, including ties. In general, on a reward-confidence round, next-edge maximization gives
\begin{equation}
\Delta_t^{\mathrm{full}}
\le 2b_t(a_t)+\lambda\bigl(r_t(a_t)-r_t(a_t^*)\bigr).
\label{eq:signedtail}
\end{equation}
If $r_t(a)\ge0$ and $\lambda\ge0$, the final term is at most $\lambda r_t(a_t)$.
\end{proposition}
\begin{proof}
A common remainder subtracts a common constant from all indices. For different remainders, compare the chosen and comparator next-edge indices, apply their reward-confidence inequalities, and subtract their omitted costs. Appendix~\ref{app:next} gives the algebra.
\end{proof}

Equation~\eqref{eq:signedtail} identifies the relevant approximation error: the difference between omitted costs, rather than their absolute size. A large common evaluator overhead need not affect selection. At a terminal child there may be no omitted execution suffix; at an internal node the next edge can be cheap while obtaining a reward remains expensive.

\begin{theorem}[Full-cost regret of next-edge UCB]
\label{thm:next}
Assume the reward observation, initialization, measurability, and separated full-cost gap conditions of Theorem~\ref{thm:log}. Maximize $I_t^{\mathrm{next}}$ with exact next-edge costs, nonnegative omitted remainders, and an integrable cumulative remainder. Then
\begin{equation}
\begin{split}
\E R_T^{\mathrm{full}}\le{}&KD\left(\frac{32\log T}{\delta^2}+6\right)\\
&+2\lambda\E\sum_{t<T}r_t(a_t).
\end{split}
\label{eq:nextbound}
\end{equation}
\end{theorem}
\begin{proof}
Proposition~\ref{prop:next} supplies the one-step error $E_t=\lambda r_t(a_t)$. Apply the counting and expectation argument of Theorem~\ref{thm:errors}. Only the selected remainder is charged, which permits the following cache accounting bound.
\end{proof}

\begin{theorem}[Cache-amortized next-edge regret]
\label{thm:nextcache}
Under Theorem~\ref{thm:next}, suppose the omitted remainder is the conditional expected value of an actual omitted cold bill $X_t$:
\begin{equation}
r_t(a_t)=\E[X_t\mid\mathcal I_t],\qquad
0\le X_t\le\operatorname{coldWork}_t.
\label{eq:tailforecast}
\end{equation}
The information $\mathcal I_t$ contains the pre-execution history and chosen action. Let actual materialization obey Theorem~\ref{thm:cache} in a fixed finite resource universe $\mathcal R$, starting from fixed $M_0$, with measurable integrable bills. Then
\begin{equation}
\begin{split}
\boxed{\E R_T^{\mathrm{full}}\le
KD\left(\frac{32\log T}{\delta^2}+6\right)
+2\lambda B_{\mathrm{cold}}},\\
B_{\mathrm{cold}}=W(\mathcal R)-W(M_0).
\end{split}
\label{eq:nextcache}
\end{equation}
For fixed $\mathcal R,K,D,\delta,\lambda$, this is $O(\log T)$ full-cost local regret.
\end{theorem}
\begin{proof}
Conditional expectation and finite-sum integration give
\[
\E\sum_{t<T}r_t(a_t)=\E\sum_{t<T}X_t
\le\E\sum_{t<T}\operatorname{coldWork}_t\le B_{\mathrm{cold}}.
\]
Substitute into~\eqref{eq:nextbound}. The last inequality is the qid--prefix conservation bound, not an assumed bound on forecasts.
\end{proof}

This result directly connects execution reuse to search allocation: finite persistent materialization can absorb the cost truncation error of next-edge UCB. The additive constant can be large, so identical asymptotic order does not imply comparable performance at practical budgets. If the resource universe grows with $T$, the displayed bound retains its $B_{\mathrm{cold}}(T)$ term and need not be logarithmic.

Repeated action-dependent proposer, repair, or scoring costs do not belong to this finite materialization budget. Cost prediction errors, mismatched completion laws, and adaptive reward drift also require the error terms of Section~\ref{sec:regret}. Equation~\eqref{eq:tailforecast} requires the actual execution law; an arbitrary uniform-terminal reference forecast does not establish it.

\begin{theorem}[Recurring omitted costs can force linear regret]
\label{thm:nextlinear}
Let rewards be deterministic with $\mu_A=0.8$, $\mu_B=0.7$, exact next-edge costs $e_A=0.4$, $e_B=0.1$, and recurring omitted costs $r_A=0$, $r_B=0.3$. Both full costs equal $0.4$. With $\lambda=1$, initialize $A$ then $B$ and subsequently maximize $I_t^{\mathrm{next}}$. Every such trajectory and $T\ge1$ satisfy
\begin{equation}
R_T^{\mathrm{full}}\ge T/10-40\log T-2/5.
\label{eq:nextlower}
\end{equation}
\end{theorem}
\begin{proof}
Next-edge net means are $(0.4,0.6)$, while full-cost net means are $(0.4,0.3)$. The former favor $B$, the latter $A$. If $R_T^{\mathrm{next}}$ is regret for the next-edge surrogate, then $R_T^{\mathrm{full}}=T/10-R_T^{\mathrm{next}}/2$. The exact-reward UCB count bound gives $R_T^{\mathrm{next}}\le80\log T+4/5$, proving the claim.
\end{proof}

Next-edge UCB can therefore have logarithmic regret for its own surrogate and linear regret for the full-cost target. Theorem~\ref{thm:nextcache} explains a structural regime in which this divergence is controlled.
